% Declaration on the use of large language models
% Many journals now require explicit disclosure of LLM usage.
% Adapt this section to accurately reflect your project's use.
% Reference policy: https://www.iscb.org/iscb-policy-statements/iscb-policy-for-acceptable-use-of-large-language-models

% [OPTIONAL] If LLMs were used in the research itself, describe that here:
In our work, as reported in the main text, LLMs were used as part of the research to assess reporting quality of MR studies.

% [OPTIONAL] If LLMs were used as coding/writing assistants, describe that here:
LLMs have been used as coding agents to assist:

\begin{itemize}
  \item for research analysis, the implementation of code (Python web crawler and annotation tools) and documentation;
  \item for the final proof-reading of the manuscript.
\end{itemize}

% coding agents
% List the coding agent tools used (e.g. GitHub Copilot, OpenCode, Cursor):
Coding agents used in our work include GitHub Copilot.
% models
% List the specific LLM models used:
The LLMs used include GPT-5 and Claude Opus 4.
% model access
% Describe how the models were accessed (API, web interface, IDE plugin, etc.):
The LLMs are accessed via API and IDE plugins.

% Describe provenance of figures and tables:
Figures and tables used in the manuscript are generated from analysis code
which AI was used in their implementation, but are not directly generated by AI.

% other tools
% List non-AI writing/proofreading tools used:
Other non-AI / non-LLM tools in proof-reading the manuscript include
textidote (\url{https://github.com/sylvainhalle/textidote}) and
typos (\url{https://github.com/crate-ci/typos}).

The final manuscript is fully drafted by the authors and has been reviewed and edited by all co-authors.
Prior to submission we have reviewed ISCB Policy for Acceptable Use of Large Language Models
(\url{https://www.iscb.org/iscb-policy-statements/iscb-policy-for-acceptable-use-of-large-language-models})
and follow its guidelines on the acceptable uses.
